\section{Introduction}
\label{sec:intro}

The rapid scaling of artificial intelligence has turned data-center electricity consumption into a material driver of global carbon emissions. Training and serving large models now account for a fast-growing share of compute demand, and the energy this requires has repeatedly been shown to carry a substantial carbon cost \citep{strubell2019energy, patterson2022carbon, faiz2024llmcarbon}. As AI workloads expand faster than electricity grids can decarbonize, the carbon emitted by the electricity a workload consumes while it runs, known as its operational carbon, has become a first-order sustainability concern and a recognized lever for aligning computing with climate-change mitigation \citep{kaack2022aligning, li2025llm}.

This carbon is partly controllable because the carbon intensity of grid electricity, the grams of CO\textsubscript{2}-equivalent emitted per kilowatt-hour, is not constant. It varies hour to hour as the generation mix shifts between fossil and renewable sources, and it varies region to region with the local grid, so that a kilowatt-hour drawn from a hydro- and nuclear-rich grid at night can be many times cleaner than one drawn from a coal-heavy grid at peak. Because a large fraction of data-center load is delay-tolerant, an operator can reduce the carbon emitted to complete the same computation by deferring flexible work to cleaner hours and routing it to cleaner regions, without buying new hardware or curtailing the workload \citep{radovanovic2023, acun2023carbonexplorer, lindberg2022geographic}.

Prior work, however, leaves two gaps. The first is methodological: much carbon-aware scheduling relies on heuristic rules that offer no guarantee of optimality and can violate operational limits, so an observed saving cannot be attributed to the approach rather than to one rule on one dataset \citep{radovanovic2023, wijayawardana2025scheduling}. The second concerns scope: the temporal and spatial dimensions are usually optimized separately rather than jointly, which leaves unresolved how much of the saving comes from when load runs versus where it runs \citep{liu2012, sukprasert2024limitations}. These gaps motivate the central research question of this thesis: \emph{how can a deterministic optimization model, using carbon intensity and the carbon-free energy fraction as scheduling signals, minimize data-center carbon emissions through joint temporal and spatial workload scheduling across multiple grid regions, subject to realistic operational constraints?}

The thesis answers this question by formalizing carbon-aware scheduling as a linear program situated within a broader research framework that translates grid carbon signals into scheduling decisions \citep{colangelo2025}. The model is driven by two hourly signals, carbon intensity and the carbon-free energy fraction, where the latter is the share of electricity supplied by carbon-free sources. Both signals are used because two hours with identical carbon intensity can differ in how structurally clean they are, and using both lets the model favour electricity that is clean because of renewable supply rather than electricity that is only incidentally low-carbon \citep{riepin2024cfe247costs, chien2023reducing}. The model is solved exactly, optimizes both dimensions jointly, and is evaluated on a two-year backtest across five structurally diverse grid regions.

On this basis the thesis makes three main contributions. First, it develops a provably optimal joint spatio-temporal scheduling model that dispatches flexible load against real hourly grid carbon-intensity and carbon-free-energy signals, with the latter entering the objective as a structural-cleanliness discount; this distinguishes it both from heuristic schedulers and from the closest exact formulation, which optimizes dedicated clean-energy procurement rather than dispatch against grid signals. Second, it identifies geographic flexibility as the dominant determinant of achievable carbon savings, exceeding every other modelling parameter by an order of magnitude in a one-at-a-time sensitivity analysis. Third, it decomposes the realized saving into spatial and temporal components and shows that routing load to cleaner regions accounts for almost all of it, whereas shifting load in time alone can increase carbon, thereby clarifying which form of flexibility actually drives the result.

The remainder of the thesis is organized as follows. Section~\ref{sec:litreview} reviews the literature on carbon-aware computing and positions the model against prior heuristic and optimization-based work. Section~\ref{sec:problem} describes the carbon data, the five target regions, and the exploratory analysis that motivates the modelling choices. Section~\ref{sec:formulation} presents the optimization model, its constraints, and its implementation. Section~\ref{sec:results} reports the computational results, and Section~\ref{sec:conclusion} concludes with the study's limitations and directions for future work.
